\documentclass[conference]{IEEEtran}
\usepackage[utf8]{inputenc}
\usepackage[T1]{fontenc}
\usepackage{cite}
\usepackage{graphicx}
\usepackage{amsmath}
\usepackage{booktabs}
\usepackage{hyperref}

\begin{document}

\title{UES-Br and LIBRAS-Mediated Music Learning: Exploratory Qualitative Study
of Deaf Adolescents' Interaction with LIBRAS Academy}

\author{
  \IEEEauthorblockN{%
    Guilherme Henrique de Oliveira Cestari\IEEEauthorrefmark{1},
    Ismar Frango Silveira\IEEEauthorrefmark{2},
    Cristiano da Silva Benites\IEEEauthorrefmark{3},
    Jefferson de Oliveira Silva\IEEEauthorrefmark{1}%
  }
  \IEEEauthorblockA{%
    \IEEEauthorrefmark{1}Pontifical Catholic University of São Paulo, São Paulo, Brazil\\
    \IEEEauthorrefmark{2}Mackenzie Presbyterian University, São Paulo, Brazil\\
    \IEEEauthorrefmark{3}Institute of Technology and Leadership, São Paulo, Brazil\\
    ghocestari@pucsp.br%
  }
}

\maketitle

\begin{abstract}
This paper reports a qualitative exploratory study of how six deaf adolescents, aged
13 to 15, interacted with LIBRAS Academy, a digital interface for music education
mediated by Brazilian Sign Language (LIBRAS). The study was conducted on April 6,
2026, at a specialized institutional site for deaf education and hearing, voice, and
language support. The empirical corpus comprises a semi-structured first-contact testing
protocol, observational annotations from individual sessions, open-ended comments, and
responses to the Brazilian Portuguese adaptation of the User Engagement Scale (UES-Br).
The analysis characterizes interaction patterns, interface barriers, mediation demands,
and the limits of interpreting UES-Br responses in this context. The sessions showed
that participants could complete the core activities and displayed different forms of
engagement, including
reflective exploration, efficient task completion, self-correction, selective
persistence, and image-based progression under reading constraints. Barriers emerged
around task onboarding, visual layout, game pace, scroll mechanics, recognition
reliability, and abstract questionnaire vocabulary. Correcting negatively worded
UES-Br items was essential for interpretation, but the questionnaire remained
dependent on mediation and vulnerable to semantic difficulty and possible
acquiescence. The paper contributes empirical evidence about LIBRAS-mediated music
learning and a methodological reflection on using UES-Br with deaf adolescents.
\end{abstract}

\begin{IEEEkeywords}
Deaf adolescents; LIBRAS; music education; UES-Br; inclusive learning technologies;
accessibility evaluation.
\end{IEEEkeywords}

\section{Introduction}

For many deaf students, music education is still organized around sounds that are
expected to be heard, named, reproduced, and evaluated through hearing-centered
pedagogical routines. Yet musical experience can also be approached through visual
patterns, rhythm, gesture, color, memory, embodied movement, and shared symbolic
activity. Digital learning environments can make these forms of access more concrete,
but only when they are designed and evaluated from the communicational repertoires of
deaf learners. The core difficulty is therefore not simply the absence of music software
for deaf students. It is the lack of evidence about how deaf students interpret
instructions, negotiate support, and remain engaged with such interfaces when they
encounter them in educational settings.

This problem has methodological consequences. If accessible educational interfaces are
evaluated mainly through design intentions, technical demonstrations, or generic usability
assumptions, researchers may overestimate their accessibility. The presence of sign
language, visual elements, or gesture recognition does not guarantee that instructions
are clear, that task rhythm fits the user, that visual layout supports interpretation, or
that a post-use questionnaire is linguistically accessible. In this context, engagement is
also difficult to interpret. A participant may complete tasks efficiently without wanting
to continue, persist in one activity while rejecting another, or report positive
impressions while requiring extensive mediation. For inclusive learning technologies,
these differences matter because they affect not only how users respond to the platform,
but also the reliability of the evidence used to improve it.

Relevant prior work addresses parts of this challenge. Benites' research on music
education for deaf children proposed and developed LIBRAS Academy, an interface that
combines musical notes, colors, Brazilian Sign Language (LIBRAS), and
artificial-intelligence-supported gesture interaction~\cite{benites2024,benites2024laclo}.
Research on assistive educational technology grounded in semiotic theory has shown how
images, written language, and LIBRAS signs can be combined in learning software for deaf
and hearing children~\cite{rocha2016,rocha2023uais}. In Human-Computer Interaction (HCI),
the User Engagement Scale (UES) and its Brazilian Portuguese adaptation, UES-Br, provide
a structured way to describe perceived engagement across dimensions such as focused
attention, perceived usability, aesthetic appeal, and
reward~\cite{obrien2010ues,miranda2021uesbr}. These contributions leave open an empirical
question: how do deaf adolescent users interact with a LIBRAS-mediated music learning
interface for the first time, and how should engagement instrument responses be
interpreted when the questionnaire requires translation and adult support?

The research gap addressed in this paper is twofold. Empirically, there is still limited
evidence, especially in the Brazilian and Latin American context, on how deaf students
interact with sign-language-mediated music learning interfaces in first-contact sessions.
Methodologically, little is known about how UES-Br responses should be interpreted when
respondents are deaf adolescents and the questionnaire requires translation, explanation,
or adult mediation. This paper does not claim to measure learning gains, validate UES-Br
for this population, or provide statistical evidence of platform effectiveness. Instead,
it offers an exploratory, qualitative account of interaction and engagement in
first-contact sessions.

The study was conducted on April 6, 2026, at the Children's Hearing Center (CeAC),
within the Division of Education and Rehabilitation of Communication Disorders (DERDIC),
at the Pontifical Catholic University of São Paulo (PUC-SP), in São Paulo, Brazil. Six
deaf adolescents aged 13 to 15 participated in individual first-contact sessions with
LIBRAS Academy, followed by a mediated UES-Br questionnaire. The analysis is guided by
three research questions:

\begin{itemize}
  \item RQ1: How do deaf adolescents aged 13 to 15 interact with LIBRAS Academy in
  individual first-contact sessions?
  \item RQ2: What forms of engagement are suggested by triangulating behavioral
  observation, UES-Br responses, and open-ended feedback?
  \item RQ3: How do mediation conditions and linguistic accessibility affect the
  interpretability of UES-Br responses in this context?
\end{itemize}

The remainder of the paper is organized as follows. Section II reviews related work on
DHH technologies, LIBRAS-mediated music learning, engagement measurement, and accessible
evaluation methods, building toward the gap the study addresses. Section III describes
the study context, participants, protocol, UES-Br scoring, and qualitative analysis
procedure. Section IV reports the results, beginning with observed interaction patterns
and interface barriers before treating UES-Br as complementary evidence. Section V
discusses the empirical and methodological implications of the findings, including
limitations and directions for future work. Section VI concludes the paper.

\section{Related Work}

Music education for deaf learners has long been designed around hearing-centered
assumptions, yet deaf individuals can perceive rhythm, vibration, and visual patterns in
music, and digital interfaces can extend those channels of access. Benites addressed this
design challenge through a research program that produced LIBRAS Academy, a music
learning environment in which notes are associated with colors, handshapes, and
gesture-recognition feedback mediated by Brazilian Sign Language
(LIBRAS)~\cite{benites2024,benites2024laclo}. In this approach, musical notes are not
only heard; they are seen, enacted, named, and compared through sign-based interaction,
allowing a learner without residual hearing to navigate the interface through visual and
LIBRAS-mediated cues. Rocha and colleagues developed complementary principles for
LIBRAS-based educational software grounded in semiotic theory~\cite{rocha2016,rocha2023uais},
showing that when images, written language, and signs are combined in educational
software, their value depends on whether learners can connect those resources during
activity, not on whether each is individually present. Applied to LIBRAS Academy, this
principle means that a musical note, a color, a handshape, a written label, and a
feedback message form a network: accessibility depends on the coherence of the whole
arrangement, not only on the presence of LIBRAS content. Neither Benites nor Rocha,
however, evaluated how deaf adolescents actually interact with such an interface under
realistic first-contact conditions -- a gap this study addresses directly.

The broader HCI literature on digital technologies for deaf and hard-of-hearing (DHH)
learners reinforces why that gap matters. A recent systematic review of over a hundred
papers on DHH technologies identified language learning, communication support, and
family sign language practice as the dominant research areas; music education and
engagement-based interfaces for deaf learners remain underrepresented~\cite{zhao2025dhhchildren}.
Sign-language applications that do exist show that their value cannot be assumed from
design intentions: Chuan and Guardino~\cite{chuan2016smartsignplay} demonstrated that
even well-designed ASL applications for DHH children and families depend on whether
users can coordinate signs, feedback, visual prompts, and task goals in context -- an
empirical question, not a design guarantee. More fundamentally, Korte et al.
~\cite{korte2014requirements} showed that requirements and usability judgments for deaf
children cannot be inferred from adult or hearing-user expectations: language
development, communication mediation, and ways of expressing difficulty all shape how
children interpret and respond to a system. These findings generalize to post-use
evaluation: instruments and protocols designed for hearing adults may produce unreliable
or uninterpretable responses with DHH populations when applied without adaptation.

Measuring engagement provides a structured way to complement behavioral observation, but
raises its own challenges in this context. The User Engagement Scale (UES), developed by
O'Brien and Toms~\cite{obrien2010ues}, operationalizes engagement as a multidimensional
construct covering Focused Attention, Perceived Usability, Aesthetic Appeal, and Reward.
These dimensions capture aspects of perceived experience that behavioral observation
alone cannot: whether a user felt in control, found the interface visually appealing, or
considered the activity worthwhile. The UES-Br adaptation by Miranda, Li, and
Darin~\cite{miranda2021uesbr} provides a validated Portuguese-language version for
Brazilian HCI contexts, and was selected here because using an English-language
instrument with Brazilian adolescents would introduce an avoidable validity threat.
However, UES-Br was validated with adult Brazilian participants; it has not been applied
with deaf populations of any age, and several of its items -- such as ``I let myself be
carried away'' and ``gratifying'' -- involve idiomatic or abstract vocabulary that may be
difficult regardless of the language of administration. Doherty and Doherty argue that
engagement research should use a diverse methodological toolkit rather than treating any
instrument as a direct window onto the phenomenon itself~\cite{doherty2018engagement}. In
this study, that argument shapes the analysis: UES-Br responses are treated as evidence
of perceived experience, while observational records are treated as evidence of behavior,
and their convergence or divergence is analytically informative precisely because the two
do not measure the same thing.

The methodological challenge of eliciting reliable self-report from deaf adolescents runs
through all of the above. Young learners' judgments of their own experience are shaped by
affect, social context, available vocabulary, and the conditions under which a question
is asked, not only by the interaction itself~\cite{rocha2023metacog,vanloon2021}. For DHH
learners, these conditions are more demanding: spoken think-aloud, self-administered
questionnaires in written form, and post-task interviews all assume communicational
conditions that may not hold for participants responding through sign language, limited
written Portuguese, or adult mediation~\cite{albalhareth2022}. There is no established
protocol for conducting accessible UX evaluation with deaf adolescents in
sign-language-mediated settings. The present study contributes to filling that gap by
combining behavioral observation, a mediated UES-Br administration, and open-ended
comments, and by reporting not only what participants did and said, but the specific
conditions under which each form of evidence was produced.

\section{Methodology}

\textbf{Study design.} This is a qualitative exploratory study that characterizes how six
deaf adolescents interacted with LIBRAS Academy during individual first-contact sessions,
and how their perceived engagement, reported through UES-Br, can be triangulated with
behavioral observation. The design prioritizes rich description of situated interaction
over measurement of outcomes or generalization.

\textbf{Setting.} All sessions were conducted on April 6, 2026, at the Children's
Hearing Center (CeAC), within DERDIC, at PUC-SP, in São Paulo, Brazil. DERDIC is a
specialized educational and clinical center that operates at the intersection of deaf
education, hearing and language rehabilitation, LIBRAS support, and family-facing
services. This setting was chosen because it provides access to deaf adolescents in a
context where sign language use, family mediation, and inclusive education practices are
part of the institutional routine, making it ecologically appropriate for evaluating an
accessible learning interface.

\textbf{Participants.} Six deaf adolescents aged 13 to 15 were recruited through the
institution. Inclusion criteria were: being a student or service user of DERDIC/CeAC,
having at least some familiarity with LIBRAS, and having a responsible adult present
during the session. Participants had not previously used LIBRAS Academy. They are
identified by codes P1 to P6 to preserve anonymity. Table~\ref{tab:participants}
summarizes their profiles; differences in hearing history, LIBRAS background, prior music
experience, and written Portuguese fluency are analytically relevant because they shaped
how each participant interpreted instructions, managed task demands, and responded to the
post-use questionnaire.

\begin{table*}[t]
  \caption{Participant profiles.}
  \label{tab:participants}
  \footnotesize
  \begin{tabular}{@{}p{0.04\linewidth} p{0.12\linewidth} p{0.22\linewidth} p{0.20\linewidth} p{0.30\linewidth}@{}}
    \toprule
    \textbf{Code} & \textbf{Age / year} & \textbf{Hearing profile} &
      \textbf{LIBRAS} & \textbf{Music / Portuguese} \\
    \midrule
    P1 & 14; 1st HS & Moderate hearing loss &
      Since 2015 (11 yrs) &
      No music study; standard written Portuguese fluency. \\
    P2 & 14; 1st HS & Complete deafness; no variation &
      Learned in deaf family (parents deaf) &
      No music study; standard written Portuguese fluency. \\
    P3 & 14; 9th gr. & Asymmetric (right 15\%, left none) &
      4--5 years; deaf cousins &
      Some prior music exposure; standard fluency. \\
    P4 & 13; 9th gr. & Complete deafness; prior hearing variation &
      Since childhood (5 yrs) &
      About 1 year of music study; standard fluency. \\
    P5 & 15; 9th gr. & Partial hearing (left ear) &
      3 years; deaf uncle &
      No formal music study; limited written Portuguese. \\
    P6 & 15; 1st HS & Partial hearing (right ear) &
      11 years; only deaf in family &
      Likes music; stronger written Portuguese fluency. \\
    \bottomrule
  \end{tabular}
\end{table*}

\textbf{Session protocol.} Each session was conducted individually and followed a
structured script. The research team comprised three members: one researcher dedicated
to interacting with the participant and guiding the task sequence, one researcher
responsible for observational annotation and audiovisual documentation, and a LIBRAS
interpreter who translated task instructions and questionnaire items when needed. A
responsible adult was present in all sessions. Sessions began with a profile interview
answered by the responsible adult and, when possible, by the participant, covering
hearing history, LIBRAS background, prior music and technology experience. The main task
block comprised seven activities presented in fixed order: free exploration of the home
screen; ``Alphabet in LIBRAS''; ``What Are the Musical Notes?''; ``Conductor of Notes'';
``Musical Rainbow''; ``Note Collector''; and a free-choice final task. Participants could
request help at any moment and could indicate readiness to advance to the next activity.
Total session duration ranged from 7 to 18 minutes of active task time across participants.

\textbf{Data collection.} The empirical corpus comprises: manual observational
annotations recorded in real time by the documentation researcher; task progression notes
indicating which phases were reached in each activity; mediation records noting when and
how support was provided; and post-use responses to UES-Br administered through Google
Forms. The UES-Br was completed by each participant immediately after the tasks, with
verbal item-by-item translation by the LIBRAS interpreter and, when necessary, additional
clarification by the responsible adult. Participants responded at their own pace;
responses were entered into Google Forms by the participant or, when needed, by the
interpreter. Screen recordings were collected as part of the documentation protocol but
are not analyzed here as a coded corpus; observational notes are the primary qualitative
source.

\textbf{UES-Br scoring.} The instrument contains 30 Likert items (scale 1--5) grouped
into four dimensions: Focused Attention (FA, 7 items), Perceived Usability (PU, 8
items), Aesthetic Appeal (AE, 5 items), and Reward (RW, 10 items), plus one open-ended
prompt. FA and AE contain only positively worded items and were computed as item means.
PU contains seven negatively worded items (PU1--PU6 and PU8) and one positively worded
item (PU7); it is reported as both a raw mean and a corrected mean, where negatively
worded items are inverted using $\text{score}_{\text{corr}} = 6 - \text{score}_{\text{raw}}$
before averaging. RW3 is negatively worded and RW is likewise reported in both raw and
corrected form. All individual item responses were verified by manual inspection against
the original PDF exports of the Google Forms records prior to computing means. Given the
small sample and the mediated administration, UES-Br results are used descriptively and
as complementary evidence of perceived experience, not for inferential claims.

\textbf{Qualitative analysis.} Observational annotations were coded using a hybrid
thematic procedure. An initial set of codes was derived deductively from the three
research questions; additional codes were developed inductively through close reading of
the annotation corpus. The coding was performed by one researcher and reviewed in a
collaborative session with the full research team, where codes were consolidated and
discrepancies resolved by consensus. The final codebook comprised five categories:
interface barriers (layout, pace, recognition, scale); persistence and voluntary
continuation; curiosity and exploratory behavior; instrument accessibility problems; and
mediation and social context. Coded events were organized by participant and activity to
support cross-case comparison. The analysis does not reach theoretical saturation; it
identifies recurring patterns and contrastive cases across six first-contact sessions. In
the triangulation, UES-Br responses and observational codes are treated as distinct forms
of evidence -- perceived experience and observed behavior, respectively -- whose
convergence or divergence is analytically informative.

\textbf{Ethical considerations.} The study was conducted in accordance with applicable
ethical guidelines for research involving minors. Informed consent was obtained from
responsible adults; participants provided verbal assent. \textit{[Ethics approval number
to be inserted before submission.]} Participation was voluntary and participants could
withdraw at any time without consequence. Anonymity is preserved through participant
codes P1--P6; no names, school identifiers, or other direct identifiers appear in the
analysis or reporting.

\section{Results}

\subsection{Participant Context and Session Overview}

All six participants completed the core sequence of LIBRAS Academy activities in a
single first-contact session. Session duration varied from 7 minutes and 25 seconds
(P4) to 18 minutes and 9 seconds of practice time (P5), reflecting differences in pace,
reading fluency, curiosity, and need for mediation, not a simple gradient of
success or failure. The participants also differed in prior music experience: P1, P2,
and P5 reported no previous music study, P3 and P4 reported some prior exposure, and
P6 reported liking music.

Table~\ref{tab:activity-matrix} summarizes the dominant form of support observed
across the activity sequence. The codes are descriptive rather than scores:
\(\bullet\) indicates predominantly autonomous progression, \(\triangle\) indicates
local mediation at task onset or at specific moments, and \(\times\) indicates a
structural barrier that interfered with the activity, such as layout, scroll, pace, or
recognition problems.

\begin{table*}[t]
  \caption{Dominant interaction pattern by participant and activity.}
  \label{tab:activity-matrix}
  \centering
  \footnotesize
  \begin{tabular}{@{}lccccccc p{0.30\linewidth}@{}}
    \toprule
    Participant & Explore & Alphabet & Notes & Conductor & Rainbow & Collector & Free choice & Main observed issue \\
    \midrule
    P1 & \(\bullet\) & \(\bullet\) & \(\triangle\) & \(\triangle\) & \(\triangle\) & \(\bullet\) & \(\bullet\) & Needed initial mediation and zoom adjustment; questionnaire language was harder than the tasks. \\
    P2 & \(\bullet\) & \(\triangle\) & \(\bullet\) & \(\triangle\) & \(\times\) & \(\triangle\) & \(\triangle\) & Layout and game pace affected interpretation; voluntary continuation was limited. \\
    P3 & \(\bullet\) & \(\triangle\) & \(\bullet\) & \(\bullet\) & \(\bullet\) & \(\bullet\) & \(\bullet\) & Mostly fluent progression, with self-correction during free-choice play. \\
    P4 & \(\bullet\) & \(\triangle\) & \(\bullet\) & \(\bullet\) & \(\triangle\) & \(\bullet\) & \(\bullet\) & Rapid interaction and impatience with slower rhythms; selective persistence in Collector. \\
    P5 & \(\triangle\) & \(\triangle\) & \(\triangle\) & \(\bullet\) & \(\times\) & \(\bullet\) & \(\triangle\) & Reading difficulty, image-based orientation, scroll problem, and possible acquiescence in evaluation. \\
    P6 & \(\bullet\) & \(\times\) & \(\triangle\) & \(\bullet\) & \(\bullet\) & \(\bullet\) & \(\bullet\) & Strong written Portuguese; recognition problem with letter O and brief difficulty with Sol. \\
    \bottomrule
  \end{tabular}
\end{table*}

\subsection{Observed Interaction Patterns}

The first pattern was reflective exploration. P1 asked whether the tool was from Brazil
or the United States, explored letters from his name, questioned the polysemy of
``Sol,'' and turned his attention to the computer itself by commenting on the notebook's
speed and entering system settings. This behavior shows that engagement includes inquiry
about the cultural, linguistic, and technical conditions of the activity, not only task
completion.

The second pattern was practical efficiency. P2 usually understood the goal of the
activities quickly and advanced with little disruption. He remembered Do, Re, and Mi and
then identified Fa and Sol. However, efficiency did not always lead to voluntary
continuation: P2 did not initially want to continue in the final free-choice moment and
proceeded only after encouragement.

The third pattern was autonomous interpretation supported by written Portuguese fluency.
P6 showed the clearest ease with written Portuguese, understood the alphabet and
Collector activities rapidly, and generally progressed without extensive mediation. His
main difficulties were localized: a system recognition problem with the letter O and a
brief hesitation around Sol. This case matters because the most visible
barrier was partly produced by the system, not by the participant's understanding.

The fourth pattern was self-monitoring and correction. P3 showed the clearest case:
during the memory game, she initially made an error, perceived it, and corrected it.
She also continued in ``Build the Word'' with good performance. One clear indicator of
engagement in these sessions was the participant's ability to notice, repair, and
continue, not only to report positive impressions afterward.

The fifth pattern was selective persistence. P4 was the fastest participant and
sometimes acted impatiently, especially in slower or more repetitive tasks. Yet she
continued playing ``Note Collector'' until the progress bar was completed. P5 also
persisted in the practice sequence for the longest duration, despite reading difficulty
and a documented impression that he might have been trying to please the researchers.
These cases caution against equating fast completion, slow completion, or high ratings
with engagement in a direct way.

\subsection{Interface and Mediation Barriers}

The most recurrent barriers involved visual design, language, and mediation. In the
musical tasks, participants often needed support at
the beginning of an activity, but this support was usually local: once the rule or
expected action was made clear, they could continue. This was visible in P1's initial
support in the notes and conductor tasks, P2's and P4's alphabet difficulties on
specific letters, and P3's local support for particular alphabet items.

Several barriers were not about music knowledge itself. P2 had initial difficulty in
``Musical Rainbow'' because notes were not aligned on the same line. P5 encountered a
scroll/touchpad problem that made the ``see again'' button inaccessible. P6's alphabet
task was affected by a recognition failure for the letter O, which is better interpreted
as a system limitation than as a participant difficulty. P1 required a zoom adjustment,
showing that visual scale can affect access even when content is conceptually clear.

Pace also shaped access. P2 perceived ``Note Collector'' as slow and tried to click to
accelerate it. P4 understood ``Conductor of Notes'' quickly, but the mediator asked her
to slow down; in ``Musical Rainbow,'' speed and task structure appeared to interact with
her difficulties. These cases show that sign-language-mediated music learning requires
temporal rhythm, feedback, and control, not only LIBRAS videos or visual musical content.

\subsection{UES-Br as Complementary Evidence}

\begin{table*}[t]
  \caption{Descriptive UES-Br means by participant. PU and RW are reported as raw means
  and as corrected means after inverting negatively worded items ($6 - \text{score}$).}
  \label{tab:uesbr}
  \begin{tabular}{lcccccc}
    \toprule
    \textbf{Code} & \textbf{FA} & \textbf{PU raw} & \textbf{PU corr.} &
      \textbf{AE} & \textbf{RW raw} & \textbf{RW corr.} \\
    \midrule
    P1 & 3.86 & 3.00 & 3.25 & 3.00 & 3.50 & 3.50 \\
    P2 & 3.00 & 3.12 & 3.12 & 3.80 & 4.00 & 3.80 \\
    P3 & 3.43 & 3.62 & 2.62 & 3.80 & 4.20 & 4.00 \\
    P4 & 2.29 & 2.12 & 4.12 & 1.80 & 2.90 & 2.70 \\
    P5 & 3.43 & 1.62 & 4.88 & 5.00 & 4.50 & 4.90 \\
    P6 & 3.29 & 3.38 & 2.88 & 3.60 & 3.30 & 3.30 \\
    \bottomrule
  \end{tabular}
\end{table*}

Table~\ref{tab:uesbr} presents the UES-Br results, which are best read as perceived
experience under mediated conditions, not as a stand-alone measurement of engagement. They help
identify contrasts among cases, but they do not replace the observational evidence.
Some of the clearest findings were divergences between observed interaction and
questionnaire response.

The corrected PU scores changed the interpretation of several cases. Most PU items are
negative statements about difficulty, confusion, or frustration; therefore, low raw
agreement can indicate a more positive usability perception after inversion. P5's PU
increased from a raw mean of 1.62 to a corrected mean of 4.88 after inverting negatively
worded items, while P3 and P6 moved in the opposite direction, from raw PU means of 3.62
and 3.38 to corrected means of 2.62 and 2.88. These shifts illustrate why reporting raw
PU alone would be misleading. RW also required correction because RW3 is negatively worded.

P3 showed the clearest convergence between observed fluency and positive self-report,
especially in reward. P4 showed lower focused attention and aesthetic appeal, but the
observational record still documented selective persistence in ``Note Collector.'' P1
completed the tasks with local mediation but struggled with the questionnaire itself,
requiring multiple explanations for expressions such as ``I let myself be carried away,''
``aesthetically pleasing,'' and ``gratifying.'' P5 is analytically the most complex
case: he produced the highest corrected PU, AE, and RW profile, but the observer noted
a possible desire to please the researchers. For that reason, P5 should be treated
as evidence about the conditions of mediated self-report, not as simple proof of high
satisfaction.

The two open responses support this cautious reading. P2 wrote that the experience
showed the importance of trying a game in LIBRAS and created in LIBRAS, while P6 wrote
that it was good. These comments are brief, but they reinforce the relevance of
LIBRAS-mediated design while also showing the limited explanatory depth of post-task
written feedback in this setting.

Across cases, the main contrasts were not reducible to high or low engagement. P1 showed
successful task interaction but substantial questionnaire barriers; P2 combined efficient
completion with limited voluntary continuation; P3 showed the clearest convergence between
self-correction and positive reward responses; P4 combined low attention and aesthetic
scores with persistence in a progress-based task; P5 showed the sharpest divergence
between very high scores and possible acquiescence; and P6 combined autonomous interaction
with localized system-related difficulty.

\section{Discussion}

In these sessions, no single feature explained accessibility by itself. Participants moved
among LIBRAS signs, written Portuguese, musical note names, colors, hand gestures, camera
feedback, progress indicators, and mediator explanations. This extends Benites' work on
AI-assisted music education for deaf learners~\cite{benites2024,benites2024laclo} by showing
how first-contact use is negotiated across these resources, and it connects the
observations to semiotic approaches to accessible educational software~\cite{rocha2016,rocha2023uais}.

From this perspective, engagement is better read through what participants did with the
available resources than through time, task success, or questionnaire scores alone. The
same interface invited curiosity, fast completion, repair after error, selective
persistence, and image-based navigation in different cases. The divergences between
observation and UES-Br therefore matter because they show how each response was produced,
not merely whether the participant ``liked'' the interface.

The institutional setting strengthens the relevance of these findings while also
delimiting them. DERDIC and CeAC operate at the intersection of deaf education,
accessibility, hearing and language support, and family-centered services. This made the
site appropriate for observing an accessible educational technology with adolescents whose
interaction is shaped by communicational and educational histories. At the same time, the
setting means that the sessions were not neutral laboratory encounters: interpretation,
mediation, institutional familiarity, and researcher presence were all part of the
interaction.

The barriers observed in the sessions can be interpreted as breakdowns in
coordination rather than as isolated usability defects. A layout problem is not only a
visual issue if it disrupts the mapping between note sequence, color, and action; a pace
problem is not only a timing issue if it changes whether the learner can anticipate,
repeat, or control the task; a recognition error is not only a technical issue if it makes
the participant's correct gesture appear incorrect. This interpretation shifts the design
question from ``is LIBRAS present?'' to ``does the interface make the relation among
LIBRAS, written language, musical representation, gesture, and feedback legible enough
for autonomous action?''

The UES-Br findings add a methodological contribution. The instrument was useful
for structuring post-use reflection, but its scores became interpretable only when the
scoring direction, mediation process, and observed behavior were considered together.
The issue is not that UES-Br is unsuitable; rather, its use with deaf adolescents requires
attention to the language in which items are understood, the social setting in which
answers are produced, and the possibility that agreement with an item may reflect
politeness, uncertainty, or translation effects as well as perceived experience. In this
study, UES-Br is best read as complementary evidence under mediated conditions, not
as an independent validation of engagement.

\subsection{Limitations}

The study has several methodological limitations. The sample includes six participants,
all tested in a single first-contact session at the same institution. The
analysis is based on observational notes, task records, open comments, and UES-Br
responses, rather than a fully coded audiovisual corpus or repeated classroom use. The
results therefore cannot be generalized statistically and should be understood as
exploratory evidence for iterative design and future research.

A specific limitation concerns the application of UES-Br in this context. The instrument
is useful because it offers a validated Brazilian Portuguese adaptation of the User
Engagement Scale~\cite{miranda2021uesbr}. However, the observations show that some participants found portions of the questionnaire
difficult. P1
required multiple explanations, did not understand expressions such as ``I let myself be
carried away,'' and had doubts about ``aesthetically pleasing'' and ``gratifying.'' P3
also did not understand ``gratifying.'' These examples suggest that some UES-Br vocabulary
and abstraction levels may be insufficiently accessible for deaf adolescents aged 13 to
15 when responses are produced through mediated communication.

Questionnaire length may also have contributed to fatigue or reduced response precision.
After several interface tasks, a long Likert-scale instrument can become cognitively
demanding, particularly when each item may require translation, contextual explanation, or
negotiation of meaning. These factors may have affected the reliability or interpretability
of some engagement responses. This does not invalidate UES-Br; rather, it indicates that
additional adaptation is needed when the instrument is used with deaf
adolescents in mediated settings.

Another limitation concerns social desirability and acquiescence. P5's case shows why
high scores cannot be treated as direct evidence of high satisfaction without contextual
checking. The observer explicitly noted a possible desire to please the researchers,
while the same session also documented reading difficulty and image-based orientation. This
does not justify excluding P5; instead, it limits any aggregate interpretation and supports
case-based triangulation.

\subsection{Future Work}

Future work should begin with additional testing involving other students from the same
institution. This would allow the research team to examine whether the patterns observed
here recur among adolescents with different LIBRAS histories, hearing profiles, prior music
experiences, and familiarity with digital technologies. A broader participant set would
also support more systematic comparison between behavioral observation and engagement
measures.

Longitudinal studies are also needed. First-contact engagement does not necessarily predict
sustained educational value. Repeated sessions could examine whether participants become more
autonomous, whether task comprehension improves, whether musical concepts are retained,
and whether sign-language-mediated musical activities contribute to literacy, communication,
memory, or interdisciplinary learning practices. Such studies would be especially useful for
clarifying the platform's relevance beyond music learning alone.

Interface refinements should focus on the issues identified in the observations: clearer
task onboarding, adjustable visual scale, more consistent spatial arrangement of notes,
transparent camera-state indicators, and pacing options for activities that feel too slow
or too fast. Engagement evaluation methods should also be adapted for deaf adolescents,
possibly through shorter scales, simplified vocabulary, LIBRAS video items, visual anchors,
and interview prompts designed for mediated communication. A future version of this study
should also predefine how screen recordings will be coded, so that observational notes and
audiovisual evidence can be compared more systematically.

\section{Conclusion}

This paper presented a qualitative exploratory analysis of six deaf adolescents'
first-contact interaction with LIBRAS Academy during tests conducted on April 6, 2026,
at CeAC/DERDIC/PUC-SP. Rather than treating the sessions as a psychometric evaluation
or as evidence of learning gains, the study examined interaction, accessibility
barriers, and perceived experience through observational evidence, open comments, and
UES-Br responses.

The results indicate that the participants completed the proposed tasks and often
demonstrated curiosity, autonomy, self-correction, or persistence. More than that, they
show that engagement and accessibility emerged from the relation among visual design,
LIBRAS mediation, written language, musical representation, task rhythm, feedback, prior
repertoire, and the social conditions of evaluation.

For evaluation, UES-Br was useful as a structured self-report instrument, but its results
required careful scoring and triangulation. Correcting negatively worded PU and RW items
changed the interpretation of several cases, and semantic difficulty or possible
acquiescence limited the strength of questionnaire-based claims.

The main contribution of this paper is twofold: an account of how deaf adolescents
interacted with a LIBRAS-mediated music learning interface in a specialized institutional
setting, and a methodological argument for interpreting engagement questionnaires through
the communication conditions under which responses are produced.
Broader and longitudinal studies are still needed to examine learning, retention, and
sustained educational value.

\bibliographystyle{IEEEtran}
\bibliography{references}

\end{document}
